%% Question bank chapter file - generated by bank_tools.py build (do not edit while the bank is still being built:
%% it is regenerated from the per-test data in Question Bank\data). Every question is one bankq block:
%%   \begin{bankq}[id=..., source={...}, paper=.., qno=.., type=..., subtopic={...}, answer={...}, answer_src=key|solved,
%%                 working={...}, difficulty=1-10, diff_parts={K S I C F}, diff_note={...},
%%                 group=<passage id or empty>, also_in={duplicate copies in other tests}]
%% difficulty = min(10, K+S+I+C+F-1): K concepts 1-3, S steps 1-4, I insight 0-2, C computation 0-2, F format 0-1
%%   ...question LaTeX, options, figures (\bfig / \bfigi -> img/)...
%%   \end{bankq}
%% Passages shared by several questions are bankpassage blocks; their questions carry group=<passage id>.
\documentclass[10pt]{article}
\usepackage[a4paper,margin=1.8cm]{geometry}
\usepackage[T1]{fontenc}
\usepackage{amsmath,amssymb,mathtools}
\usepackage{graphicx}
\usepackage[export]{adjustbox}
\usepackage{xcolor,needspace,keyval}
\usepackage{enumitem}
\usepackage{booktabs,array,tabularx}
\usepackage[hidelinks]{hyperref}
\setlength{\parindent}{0pt}
\setlength{\parskip}{3pt}
\makeatletter
\define@key{bq}{id}{\def\bq@id{#1}}\define@key{bq}{source}{\def\bq@src{#1}}\define@key{bq}{paper}{\def\bq@paper{#1}}
\define@key{bq}{qno}{\def\bq@qno{#1}}\define@key{bq}{type}{\def\bq@type{#1}}\define@key{bq}{subtopic}{\def\bq@sub{#1}}
\define@key{bq}{answer}{\def\bq@ans{#1}}\define@key{bq}{answer_src}{\def\bq@asrc{#1}}\define@key{bq}{working}{\def\bq@work{#1}}
\define@key{bq}{group}{\def\bq@group{#1}}\define@key{bq}{also_in}{\def\bq@also{#1}}
\define@key{bq}{difficulty}{\def\bq@diff{#1}}\define@key{bq}{diff_parts}{\def\bq@dparts{#1}}\define@key{bq}{diff_note}{\def\bq@dnote{#1}}
\newenvironment{bankq}[1][]{\def\bq@work{}\def\bq@also{}\def\bq@group{}\def\bq@diff{?}\def\bq@dparts{}\def\bq@dnote{}\setkeys{bq}{#1}%
  \par\needspace{4\baselineskip}\medskip\noindent{\small\textbf{\bq@id}\hfill\color{gray}[\bq@src, Paper \bq@paper, Q\bq@qno]
  \bq@type\ $\cdot$ \bq@sub\ $\cdot$ \textbf{D\bq@diff}}\par\nopagebreak\relax}
  {\par{\small\color{teal}Answer: \textbf{\bq@ans}\ifx\bq@work\@empty\else\ (solved: \bq@work)\fi
  \ifx\bq@also\@empty\else\par\color{gray}Also in: \bq@also\fi}\par}
\define@key{bp}{id}{\def\bp@id{#1}}\define@key{bp}{source}{\def\bp@src{#1}}\define@key{bp}{questions}{\def\bp@qs{#1}}
\newenvironment{bankpassage}[1][]{\setkeys{bp}{#1}\par\needspace{4\baselineskip}\bigskip\noindent
  {\small\textbf{Passage \bp@id}\hfill\color{gray}[\bp@src; questions \bp@qs]}\par\nopagebreak\relax\itshape}{\par}
\makeatother
\newcommand{\bfig}[2]{\par\smallskip{\centering\includegraphics[width=#1pt,max width=0.62\linewidth,max height=0.32\textheight]{img/#2}\par}\smallskip}
\newcommand{\bfigi}[2]{\includegraphics[width=#1pt,max width=\linewidth,max height=0.32\textheight,valign=c]{img/#2}}
\newcommand{\bmissing}[1]{\textit{[figure missing in the source conversion: \detokenize{#1}]}}
\newenvironment{qoptions}{\begin{itemize}[leftmargin=1.8em,itemsep=1pt,topsep=2pt,label={}]}{\end{itemize}}
\newcommand{\opt}[2]{\item[(#1)] #2}
\providecommand{\fwstart}{}\providecommand{\fwend}{}\providecommand{\qfullsplit}{}
\newcolumntype{L}[1]{>{\raggedright\arraybackslash}p{#1}}
\newcolumntype{C}[1]{>{\centering\arraybackslash}p{#1}}
\newcolumntype{R}[1]{>{\raggedleft\arraybackslash}p{#1}}
\providecommand{\sectionhdr}[1]{}\providecommand{\subsectionhdr}[1]{}


\begin{document}
{\Large\bfseries Continuity \& Differentiability --- Question Bank}\par\smallskip
{\small Mathematics chapter M23. Questions from all converted test papers, grouped by source series. Answers: printed key where it matches the paper, otherwise solved (answer\_src=solved, with working). Exact repeats across papers appear once, with also\_in listing the other copies.}\par\medskip\hrule

\section*{Score-I papers}
\begin{bankq}[id=SI12-P1-MATH-S3-Q09, source={Score-I - Test 12 (09 Mar 2022)}, paper=1, qno=9, type=numerical, subtopic={Differentiability of modulus and greatest-integer functions}, answer={180}, answer_src=key, difficulty=8, diff_parts={K2 S4 I1 C1 F1}, diff_note={tan(pi[x^2])=0 so f=ax^3+b on (0,1]; continuity and derivative match at x=1 give a=1/6, b=pi/4-13/6}]
If $f(x)=\begin{cases}\dfrac{\tan[x^2]\pi}{ax^2}+ax^3+b, & 0<x\le1\\[6pt] 2\cos\pi x+\tan^{-1}x, & 1<x\le2\end{cases}$ is differentiable in $(0,2]$, then $a=\dfrac1{k_1}$ and $b=\dfrac\pi4-\dfrac{26}{k_2}$.\\
Find $k_1^2+k_2^2$ {where $[.]$ denotes greatest integer function}.
\end{bankq}
\begin{bankq}[id=SI12-P2-MATH-S3-Q10, source={Score-I - Test 12 (09 Mar 2022)}, paper=2, qno=10, type=matching, subtopic={Types of discontinuity}, answer={P$\to$1, Q$\to$1, R$\to$3}, answer_src=key, difficulty=7, diff_parts={K3 S3 I1 C0 F1}, diff_note={Continuity count of piecewise function, integration by parts for f(2), maximal integral between roots}]
\begin{center}\footnotesize
\begin{tabular}{|p{0.6cm}|p{5.3cm}|p{0.6cm}|p{1.0cm}|}
\hline
 & \centering List-I & & \centering List-II\tabularnewline
\hline
(P) & The function $f(x)$ is defined as $f(x)=\begin{cases}(-1)^{[x^2]} & \text{if } x<0\\[2pt] \displaystyle\lim_{n\to\infty}\dfrac1{1+x^n} & \text{if } x\ge0\end{cases}$ where $[x]$ denotes the greatest integer function. Number of integral values of $x$ in $[-2,9]$ for which $f(x)$ is continuous is & (1) & 9\\
\hline
(Q) & $f:R\to R$ and satisfies $f(1)=-1$, $f'(1)=3$. If $\displaystyle\int_1^2(2-x)f''(x)\,dx=7$, then $f(2)$ equals & (2) & 1\\
\hline
(R) & If $\displaystyle\int_a^b\left(4+3x-x^2\right)dx$ is maximum then $(a+b)$ is equal to \ $(b>a)$ & (3) & 3\\
\hline
\end{tabular}
\end{center}
\end{bankq}

\section*{TOAS papers}
\begin{bankq}[id=TOAS06-P1-MATH-S3-Q01, source={TOAS - Test 06 (17 Apr 2022)}, paper=1, qno=1, type=numerical, subtopic={Differentiability of modulus and greatest-integer functions}, answer={1.00}, answer_src=key, difficulty=5, diff_parts={K1 S3 I1 C0 F1}, diff_note={Double zeros at 0 and 1 are smooth; only the simple zero at x=2 is a corner}]
Number of points where \[ f(x)=\left|(2^x-1)(3^x-3)\,x(x-1)(x-2)\right| \] is not differentiable, is
\end{bankq}
\begin{bankq}[id=TOAS07-P1-MATH-S2-Q09, source={TOAS - Test 07 (01 May 2022)}, paper=1, qno=9, type=multi, subtopic={Types of discontinuity}, answer={A,B,C,D}, answer_src=key, difficulty=7, diff_parts={K2 S3 I1 C1 F1}, diff_note={Composite gofog is 1/(1/x-1)... where [.] jumps}]
Let $f(x)=x-1$ \& $g(x)=\dfrac1x$, then the set of points, where $[(gofog)(x)]$ is discontinuous is ($[.]$ denotes greatest integer function) -
\begin{qoptions}
\opt{A}{$\{0\}$}
\opt{B}{$\left\{\dfrac{n}{n+1}\right\}$, $n\in I-\{-1\}$}
\opt{C}{$\left\{\dfrac{n}{n+1}\right\}\cup\{1\}$, $n\in I-\{-1\}$}
\opt{D}{$\{0,1\}$}
\end{qoptions}
\end{bankq}
\begin{bankpassage}[id=TOAS07-P1-MATH-S3-G5-6, source={TOAS - Test 07 (01 May 2022)}, questions={5-6}]
Consider $f(x)=x^4-3x^3+2$ and $h(x)=\dfrac{f(x)}{g(x)}$ be a rational function.\\
(i) $\displaystyle\lim_{x\to\infty}h(x)=$ infinite\\
(ii) $\displaystyle\lim_{x\to1}h(x)=\dfrac52$\\
(iii) $h(x)$ is continuous $\forall\ x\in R-\{1,2\}$\\
(iv) $g(x)$ does not have any repeated factors.
\end{bankpassage}
\begin{bankq}[id=TOAS07-P1-MATH-S3-Q05, source={TOAS - Test 07 (01 May 2022)}, paper=1, qno=5, type=numerical, subtopic={Differentiability of modulus and greatest-integer functions}, answer={2.00}, answer_src=key, difficulty=8, diff_parts={K2 S4 I2 C0 F1}, diff_note={Rational function conditions determine g; points where sgn g is non differentiable}, group=TOAS07-P1-MATH-S3-G5-6]
Number of points where $P(x)=\mathrm{sgn}(g(x))$ is non differentiable in $R$ is (where $\mathrm{sgn}(x)$ denotes signum function of $x$)
\end{bankq}
\begin{bankq}[id=TOAS07-P1-MATH-S3-Q06, source={TOAS - Test 07 (01 May 2022)}, paper=1, qno=6, type=numerical, subtopic={Differentiability of modulus and greatest-integer functions}, answer={6.00}, answer_src=key, difficulty=6, diff_parts={K2 S3 I1 C0 F1}, diff_note={Count solutions of |sin(1/|x|)|=1/2}, group=TOAS07-P1-MATH-S3-G5-6]
Number of solutions of equation $|g(|x|)|=\dfrac12$ is -
\end{bankq}
\begin{bankq}[id=TOAS07-P2-MATH-S4-Q09, source={TOAS - Test 07 (01 May 2022)}, paper=2, qno=9, type=numerical, subtopic={Types of discontinuity}, answer={2}, answer_src=key, difficulty=7, diff_parts={K2 S3 I2 C0 F1}, diff_note={Sign changes of the cubic's roots and g: minimum discontinuities}]
Let $f(x)=g(x).\,\mathrm{sgn}(x^3+6x^2+11x+6)$, where $\mathrm{sgn}$ denotes signum function \& $g(x)=ax^3+x^2+x+1$, then minimum number of points of discontinuity of the function $f(x)$ is
\end{bankq}
\begin{bankq}[id=TOAS08-P1-MATH-S1-Q05, source={TOAS - Test 08 (15 May 2022)}, paper=1, qno=5, type=multi, subtopic={Types of discontinuity}, answer={A,B,C}, answer_src=key, difficulty=6, diff_parts={K2 S3 I1 C0 F1}, diff_note={Four statements on periodicity, differentiability, boundedness, IVT}]
Identify the CORRECT statements -
\begin{qoptions}
\opt{A}{Sum of two periodic functions may or maynot be a periodic function}
\opt{B}{A differentiable function in $R$ will always be continuous in $R$}
\opt{C}{A continuous function whose domain is in a closed interval is always bounded}
\opt{D}{In $[a,b]$ for a function $f(x)$, if $f(a).f(b)<0$, then there exists atleast one root of $f(x)=0$ in $(a,b)$}
\end{qoptions}
\end{bankq}
\begin{bankpassage}[id=TOAS08-P1-MATH-S2-G9-10, source={TOAS - Test 08 (15 May 2022)}, questions={9-10}]
Let $f(x)$ be a polynomial function satisfying the following conditions $\displaystyle\lim_{x\to\infty}\dfrac{f(x)}{|x|^3}=0$, $\displaystyle\lim_{x\to\infty}\left(\sqrt{f(x)}-x\right)=-1$ and $f(0)=0$ (Note : $[k]$ denotes greatest integer function less than or equal to $k$)\newline
On the basis of above information, answer the following questions :
\end{bankpassage}
\begin{bankq}[id=TOAS08-P1-MATH-S2-Q09, source={TOAS - Test 08 (15 May 2022)}, paper=1, qno=9, type=single, subtopic={Differentiability of modulus and greatest-integer functions}, answer={D}, answer_src=key, difficulty=6, diff_parts={K1 S3 I2 C1 F0}, diff_note={Polynomial from the two limits: points of discontinuity of its floor}, group=TOAS08-P1-MATH-S2-G9-10]
The number of points of discontinuity of $g(x)=[f(x)]$ in $[0,3]$ is
\begin{qoptions}
\opt{A}{2}
\opt{B}{3}
\opt{C}{4}
\opt{D}{5}
\end{qoptions}
\end{bankq}
\begin{bankq}[id=TOAS08-P1-MATH-S2-Q10, source={TOAS - Test 08 (15 May 2022)}, paper=1, qno=10, type=single, subtopic={Differentiability of modulus and greatest-integer functions}, answer={B}, answer_src=key, difficulty=6, diff_parts={K1 S3 I2 C1 F0}, diff_note={Points of non-derivability of |f(|x|)|}, group=TOAS08-P1-MATH-S2-G9-10]
The number of points of non-derivability of $h(x)=|f(|x|)|$ is -
\begin{qoptions}
\opt{A}{2}
\opt{B}{3}
\opt{C}{4}
\opt{D}{5}
\end{qoptions}
\end{bankq}
\begin{bankq}[id=TOAS08-P2-MATH-S1-Q01, source={TOAS - Test 08 (15 May 2022)}, paper=2, qno=1, type=multi, subtopic={Differentiability of modulus and greatest-integer functions}, answer={A,B,D}, answer_src=key, difficulty=7, diff_parts={K2 S4 I1 C0 F1}, diff_note={Piecewise function: differentiability at 0, pi and -pi}]
If $f(x)=\begin{cases}\cos^{-1}(\cos x)\,; & |x|<\pi\\[2pt] \pi\log_\pi|x|\,; & |x|\geqslant\pi\end{cases}$, then correct statement(s) is/are -
\begin{qoptions}
\opt{A}{$f(x)$ is not differentiable at three points}
\opt{B}{$f'(\pi^+)+f'(-\pi^-)=0$}
\opt{C}{$f'(\pi^+)=f'(-\pi^-)=0$}
\opt{D}{$f'(0^+)+f'(0^-)=0$}
\end{qoptions}
\end{bankq}
\begin{bankq}[id=TOAS08-P2-MATH-S1-Q05, source={TOAS - Test 08 (15 May 2022)}, paper=2, qno=5, type=multi, subtopic={Types of discontinuity}, answer={B,C}, answer_src=key, difficulty=6, diff_parts={K2 S3 I1 C0 F1}, diff_note={Limit function with x^{2n}: continuity and limits}]
If $f(x)=\displaystyle\lim_{n\to\infty}\dfrac{2x^{2n}+x+2}{x^{2n}-x^2+3}$, then which of the following statement(s) holds good ?
\begin{qoptions}
\opt{A}{$f(x)$ is continuous at $x=1$}
\opt{B}{$f(x)$ is discontinuous at $x=1$}
\opt{C}{$\displaystyle\lim_{x\to2}f(x)$ exists finitely}
\opt{D}{$\displaystyle\lim_{x\to2}f(x)$ does not exist}
\end{qoptions}
\end{bankq}
\begin{bankq}[id=TOAS08-P2-MATH-S3-Q01, source={TOAS - Test 08 (15 May 2022)}, paper=2, qno=1, type=numerical, subtopic={Differentiability of modulus and greatest-integer functions}, answer={3.00}, answer_src=key, difficulty=5, diff_parts={K1 S3 I1 C0 F1}, diff_note={Differentiability at x=2 gives a and b}]
Let $f(x)=\begin{cases}ax\,; & x<2\\ ax^2-bx+3\,; & x\geqslant2\end{cases}$ and `$f$' is differentiable $\forall\,x\in R$, then $(a+b)$ is equal to -
\end{bankq}
\begin{bankq}[id=TOAS08-P2-MATH-S3-Q05, source={TOAS - Test 08 (15 May 2022)}, paper=2, qno=5, type=numerical, subtopic={Differentiability of modulus and greatest-integer functions}, answer={9.00}, answer_src=key, difficulty=7, diff_parts={K2 S4 I1 C0 F1}, diff_note={Differentiability at 0 requires matching derivatives: p quadratic}]
Let $p_1$ and $p_2$ be two values of $p$ for which $f(x)=\begin{cases}\dfrac{x(\ln(1+x)+\ln(1-x))}{(\sec x-\cos x)}\,; & x\in(-1,0)\\[6pt] \left(p^2-3p-1\right)\sin x+x^2\,; & x\in[0,\infty)\end{cases}$ is differentiable at $x=0$ find $p_1^2+p_2^2$
\end{bankq}
\begin{bankq}[id=TOAS09-P1-MATH-S2-Q09, source={TOAS - Test 09 (30 May 2022)}, paper=1, qno=9, type=multi, subtopic={Types of discontinuity}, answer={A,C}, answer_src=key, difficulty=2, diff_parts={K1 S1 I0 C0 F1}, diff_note={Types of discontinuity: standard approach for this subtopic; 1 concept(s), 1 step(s)}]
Let $g(x)$ be a polynomial of degree 2018 and
\[f(x)=\begin{cases}g\!\left(\dfrac{1}{x^5}\right)e^{-\frac{1}{x^6}}, & x\neq0\\[4pt] 0, & x=0\end{cases}\]
then
\begin{qoptions}
\opt{A}{$f(x)$ is continuous at $x=0$}
\opt{B}{$f(x)$ is non-differentaible at $x=0$}
\opt{C}{$f'(0)=\displaystyle\lim_{x\to0}f(x)$}
\opt{D}{$\displaystyle\lim_{x\to0}f(x)$ does not exist}
\end{qoptions}
\end{bankq}
\begin{bankq}[id=TOAS09-P1-MATH-S2-Q10, source={TOAS - Test 09 (30 May 2022)}, paper=1, qno=10, type=multi, subtopic={Differentiability of modulus and greatest-integer functions}, answer={A,C,D}, answer_src=key, difficulty=6, diff_parts={K2 S3 I1 C0 F1}, diff_note={Differentiability of modulus and greatest-integer functions: standard approach for this subtopic; 2 concept(s), 3 step(s), needs a non-obvious insight}]
Let
\[f(x)=\begin{cases}-2 & \text{for } -3\leqslant x\leqslant0\\ x-2 & \text{for } 0<x\leqslant3\end{cases}\]
$g(x)=[f(|x|)]+[|f(x)|]$ where $[.]$ denotes the greatest integer function. Which of the following statement(s) is/are true ?
\begin{qoptions}
\opt{A}{$g(1.7)=-1$}
\opt{B}{$g(x)$ is not differentiable at any point}
\opt{C}{$g(x)$ is continuous at all non integral points}
\opt{D}{$g(x)$ is differentiable at some points}
\end{qoptions}
\end{bankq}
\begin{bankq}[id=TOAS09-P2-MATH-S1-Q05, source={TOAS - Test 09 (30 May 2022)}, paper=2, qno=5, type=single, subtopic={Differentiability of modulus and greatest-integer functions}, answer={B}, answer_src=key, difficulty=2, diff_parts={K1 S2 I0 C0 F0}, diff_note={Differentiability of modulus and greatest-integer functions: standard approach for this subtopic; 1 concept(s), 2 step(s)}]
If the function $g(x)=\begin{cases}ae^x, & x\leqslant0\\ b\cos x+x, & x>0\end{cases}$ is differentiable, then the value of $a^2+b^2$ is -
\begin{qoptions}
\opt{A}{5}
\opt{B}{2}
\opt{C}{1}
\opt{D}{13}
\end{qoptions}
\end{bankq}
\begin{bankq}[id=TOAS09-P2-MATH-S1-Q20, source={TOAS - Test 09 (30 May 2022)}, paper=2, qno=20, type=single, subtopic={Differentiability of modulus and greatest-integer functions}, answer={C}, answer_src=key, difficulty=1, diff_parts={K1 S1 I0 C0 F0}, diff_note={Differentiability of modulus and greatest-integer functions: standard approach for this subtopic; 1 concept(s), 1 step(s)}]
Number of points where the function $f(x)=\max\left(\sqrt{2x-x^2},\ 2-x\right)$ is non differentiable, is
\begin{qoptions}
\opt{A}{0}
\opt{B}{1}
\opt{C}{2}
\opt{D}{3}
\end{qoptions}
\end{bankq}
\end{document}
